\let\negmedspace\undefined
\let\negthickspace\undefined
\documentclass[journal]{IEEEtran}
\usepackage[a5paper, margin=10mm, onecolumn]{geometry}
%\usepackage{lmodern} % Ensure lmodern is loaded for pdflatex
\usepackage{tfrupee} % Include tfrupee package

\setlength{\headheight}{1cm} % Set the height of the header box
\setlength{\headsep}{0mm}     % Set the distance between the header box and the top of the text

\usepackage{cite}
\usepackage{graphicx}
\usepackage{amsmath,amssymb,amsfonts,amsthm}
\usepackage{algorithmic}
\usepackage{graphicx}
\usepackage{textcomp}
\usepackage{xcolor}
\usepackage{txfonts}
\usepackage{listings}
\usepackage{enumitem}
\usepackage{mathtools}
\usepackage{gensymb}
\usepackage{comment}
\usepackage[breaklinks=true]{hyperref}
\usepackage{tkz-euclide} 
\usepackage{listings}
\usepackage{gvv}                                        
%\def\inputGnumericTable{}                                 
\usepackage[latin1]{inputenc} 
\usetikzlibrary{arrows.meta, positioning}
\usepackage{xparse}
\usepackage{color}                                            
\usepackage{array}                                            
\usepackage{longtable}                                       
\usepackage{calc}                                             
\usepackage{multirow}
\usepackage{multicol}
\usepackage{hhline}                                           
\usepackage{ifthen}                                           
\usepackage{lscape}
\usepackage{tabularx}
\usepackage{array}
\usepackage{float}
\newtheorem{theorem}{Theorem}[section]
\newtheorem{problem}{Problem}
\newtheorem{proposition}{Proposition}[section]
\newtheorem{lemma}{Lemma}[section]
\newtheorem{corollary}[theorem]{Corollary}
\newtheorem{example}{Example}[section]
\newtheorem{definition}[problem]{Definition}
\newcommand{\BEQA}{\begin{eqnarray}}
\newcommand{\EEQA}{\end{eqnarray}}
\usepackage{float}
\theoremstyle{remark}
\usepackage{circuitikz}
\usepackage{tikz} 
\graphicspath{{figs/}}

\begin{document}
\title{2.10.51}
\author{EE25BTECH11022 - sankeerthan}
% \maketitle
% \newpage
% \bigskip
{\let\newpage\relax\maketitle}

\renewcommand{\thefigure}{\theenumi}
\renewcommand{\thetable}{\theenumi}
\setlength{\intextsep}{10pt} % Space between text and floats



\textbf{Question}:If $\mathbf{a}, \mathbf{b}, \mathbf{c}$ are three non-zero, non-coplanar vectors and \\

$\mathbf{b}_1 = \mathbf{b} - \frac{\mathbf{b} \cdot \mathbf{a}}{|\mathbf{a}|^2} \mathbf{a}, \qquad \mathbf{b}_2 = \mathbf{b} + \frac{\mathbf{b} \cdot \mathbf{a}}{|\mathbf{a}|^2} \mathbf{a}$, \\
$\begin{aligned}
\mathbf{c}_1 = \mathbf{c} - \frac{\mathbf{c} \cdot \mathbf{a}}{|\mathbf{a}|^2} \mathbf{a} + \frac{\mathbf{b} \cdot \mathbf{c}}{|\mathbf{c}|^2} \mathbf{b}_1, 
\mathbf{c}_2 = \mathbf{c} - \frac{\mathbf{c} \cdot \mathbf{a}}{|\mathbf{a}|^2} \mathbf{a} + \frac{\mathbf{b}_1 \cdot \mathbf{c}}{|\mathbf{b}_1|^2} \mathbf{b}_1, \\
\mathbf{c}_3 = \mathbf{c} - \frac{\mathbf{c} \cdot \mathbf{a}}{|\mathbf{c}|^2} \mathbf{a} + \frac{\mathbf{b} \cdot \mathbf{c}}{|\mathbf{c}|^2} \mathbf{b}_1,
\mathbf{c}_4 = \mathbf{c} - \frac{\mathbf{c} \cdot \mathbf{a}}{|\mathbf{c}|^2} \mathbf{a} = \frac{\mathbf{b} \cdot \mathbf{c}}{|\mathbf{b}|^2} \mathbf{b}_1,
\end{aligned}$ \\

then the set of orthogonal vectors is
\begin{enumerate} 
\begin{multicols}{4}
\item($\mathbf{a}, \mathbf{b}_1, \mathbf{c}_3$) 
\item($\mathbf{a}, \mathbf{b}_1, \mathbf{c}_2$)
\item($\mathbf{a}, \mathbf{b}_1, \mathbf{c}_1$)
\item($\mathbf{a}, \mathbf{b}_2, \mathbf{c}_2$) 
\end{multicols}
\end{enumerate}
\textbf{solution}:
Given three non-zero, non-coplanar vectors $\mathbf{a}$, $\mathbf{b}$, $\mathbf{c}$, define: \\
\begin{align}
\mathbf{b}_1 &= \mathbf{b} - \frac{\mathbf{b} \cdot \mathbf{a}}{|\mathbf{a}|^2}\mathbf{a} \\
\mathbf{c}_2 &= \mathbf{c} - \frac{\mathbf{c} \cdot \mathbf{a}}{|\mathbf{a}|^2}\mathbf{a} + \frac{\mathbf{b}_1\cdot\mathbf{c}}{|\mathbf{b}_1|^2}\mathbf{b}_1 
\end{align}

\begin{align}
\mathbf{a}\cdot\mathbf{b}_1 = \mathbf{a}\cdot\mathbf{b} - \frac{\mathbf{b}\cdot\mathbf{a}}{|\mathbf{a}|^2}|\mathbf{a}|^2 = 0.
\end{align}
$\mathbf{b}_1$ is orthogonal to $\mathbf{a}$.
\begin{align}
\mathbf{a}\cdot\mathbf{c}_2&= \mathbf{a}\cdot\mathbf{c}-\frac{\mathbf{c}\cdot\mathbf{a}}{|\mathbf{a}|^2}|\mathbf{a}|^2 + \frac{\mathbf{b}_1\cdot\mathbf{c}}{|\mathbf{b}_1|^2}\mathbf{a}\cdot\mathbf{b}_1 \\
&=\mathbf{a}\cdot\mathbf{c}-\mathbf{c}\cdot\mathbf{a}+0 \\
&=0
\end{align}
($\mathbf{a}\cdot\mathbf{b}_1 = 0$ as above)
\begin{align}
\mathbf{b}_1\cdot\mathbf{c}_2&= \mathbf{b}_1\cdot\mathbf{c} - \frac{\mathbf{c}\cdot\mathbf{a}}{|\mathbf{a}|^2}\mathbf{b}_1\cdot\mathbf{a} + \frac{\mathbf{b}_1\cdot\mathbf{c}}{|\mathbf{b}_1|^2}|\mathbf{b}_1|^2 \\
&=\mathbf{b}_1\cdot\mathbf{c}-\mathbf{b}_1\cdot\mathbf{c} \\
&=0
\end{align}
($\mathbf{b}_1\cdot\mathbf{a} = 0$ from above) \\
$\mathbf{c}_2$ is orthogonal to both $\mathbf{a}$ and $\mathbf{b}_1$\\
Therefore, the orthogonal set is:
($\mathbf{a},\ \mathbf{b}_1,\ \mathbf{c}_2$)
\end{document}